% Figure: a Gerber (cantilever-hinged) beam. A continuous span is made
% determinate by inserting an internal hinge, which supplies the extra
% equation that a fourth support would otherwise render indeterminate. The
% suspended span rests on the cantilever arm through the hinge; the hinge
% transmits shear but no moment, so the moment diagram passes through zero
% there.
\begin{figure}[htbp]
\centering
\begin{tikzpicture}[>=latex, x=1.0cm, y=1.0cm]
  % beam line
  \draw[engblue, very thick] (0,0) -- (8,0);
  % supports: pin at A (x=0), roller at B (x=3), roller at E (x=8)
  \coordinate (A) at (0,0);
  \coordinate (B) at (3,0);
  \coordinate (H) at (5,0);   % internal hinge
  \coordinate (E) at (8,0);
  % pin support A (triangle + ground)
  \fill[engblack] (A) circle (1.6pt);
  \draw[engblack, thick] (A) -- ++(-0.28,-0.5) -- ++(0.56,0) -- cycle;
  \draw[engblack] (-0.32,-0.5) -- (0.32,-0.5);
  \foreach \x in {-0.28,-0.14,...,0.28} { \draw[enggray, thin] (\x,-0.5) -- (\x-0.12,-0.66); }
  \node[font=\scriptsize, anchor=north east] at (A) {$A$};
  % roller support B
  \fill[engblack] (B) circle (1.6pt);
  \draw[engblack, thick] (B) -- ++(-0.28,-0.4) -- ++(0.56,0) -- cycle;
  \fill[engblack] ($(B)+(-0.16,-0.5)$) circle (1.2pt);
  \fill[engblack] ($(B)+(0.16,-0.5)$) circle (1.2pt);
  \draw[engblack] ($(B)+(-0.34,-0.62)$) -- ($(B)+(0.34,-0.62)$);
  \node[font=\scriptsize, anchor=north east] at (B) {$B$};
  % roller support E
  \fill[engblack] (E) circle (1.6pt);
  \draw[engblack, thick] (E) -- ++(-0.28,-0.4) -- ++(0.56,0) -- cycle;
  \fill[engblack] ($(E)+(-0.16,-0.5)$) circle (1.2pt);
  \fill[engblack] ($(E)+(0.16,-0.5)$) circle (1.2pt);
  \draw[engblack] ($(E)+(-0.34,-0.62)$) -- ($(E)+(0.34,-0.62)$);
  \node[font=\scriptsize, anchor=north west] at (E) {$E$};
  % internal hinge at H (drawn as a small open circle joining two pen strokes)
  \draw[engblue, very thick] (0,0.06) -- (5,0.06);
  \draw[engblue, very thick] (5,-0.06) -- (8,-0.06);
  \fill[white] (H) circle (2.4pt);
  \draw[englime!55!black, very thick] (H) circle (2.4pt);
  \node[englime!55!black, font=\scriptsize, anchor=south] at ($(H)+(0,0.14)$) {hinge $H$};
  % loads: UDL over whole span (arrows) + point load P on suspended span
  \foreach \x in {0.4,1.2,...,7.6} { \draw[->, engred, thin] (\x,0.95) -- (\x,0.12); }
  \draw[engred, thin] (0.2,0.95) -- (7.8,0.95);
  \node[engred, font=\scriptsize, anchor=south] at (1.6,0.95) {$w_0$};
  \draw[->, very thick, engred] (6.5,1.5) -- (6.5,0.12);
  \node[engred, font=\scriptsize, anchor=south] at (6.5,1.5) {$P$};
  % dimension line
  \draw[<->, enggray, thin] (0,-1.0) -- (3,-1.0);
  \node[enggray, font=\scriptsize, anchor=north] at (1.5,-1.0) {$a$};
  \draw[<->, enggray, thin] (3,-1.0) -- (5,-1.0);
  \node[enggray, font=\scriptsize, anchor=north] at (4.0,-1.0) {$b$};
  \draw[<->, enggray, thin] (5,-1.0) -- (8,-1.0);
  \node[enggray, font=\scriptsize, anchor=north] at (6.5,-1.0) {$c$};
\end{tikzpicture}
\caption[Gerber cantilever-hinged beam]{A Gerber beam: a two-span beam on three
supports made statically determinate by an internal hinge $H$. Three supports
supply four reaction components (the pin at $A$ two, the rollers at $B$ and $E$
one each), one more than the three planar equilibrium equations of the whole
beam. The internal hinge supplies the missing equation: it transmits shear but
carries no moment, so the moment of all forces on the suspended span $HE$ about
$H$ vanishes, closing the system at four equations for four unknowns. The
suspended span $HE$ hangs from the cantilever arm $BH$ through the hinge, and
because the hinge carries no moment the bending-moment diagram passes through
zero there, the hinge placed by the designer at the point where the continuous
beam's moment would naturally cross zero so that inserting it costs no bending
strength. The Gerber scheme let long railway and roadway spans be built
determinate, each span analysable on its own once the hinge reaction is passed
across.}
\label{fig:vol03ch02:gerber-beam}
\end{figure}
